% !PRINTED 140
\weberpage{140}
%<<CONT

arrange according to powers of $x$, then the coefficients of the separate powers of $x$ are likewise symmetric functions. For they change no more under interchange of the $\alpha$'s than does the function $f(x)$ (\S\,1). We put
\begin{equation}
 f(x)=x^n+a_1x^{n-1}+a_2x^{n-2}+\cdots+a_{n-1}x+a_n.
\tag{2}
\end{equation}

The functions $a_1,a_2,\ldots,a_n$ are the coefficients of the algebraic equation whose roots are $\alpha_1,\alpha_2,\ldots,\alpha_n$. They are called the \gesperrt{symmetric fundamental functions}. We already formed these symmetric fundamental functions in \S\,7; as we recall, $a_1$ is the sum of the $\alpha$'s taken with negative sign, $a_2$ the sum of the products two at a time, $a_3$ the sum of the products three at a time taken with negative sign, and so on; finally $a_n$ is the product of all the $\alpha$'s with positive or negative sign according as $x$ is even or odd.
% [sic] The registered source prints x here; the parity determining the sign is that of the degree n.
We write, for short,
\begin{equation}
\begin{aligned}
 a_1&=-\sum \alpha_1,\\
 a_2&=+\sum \alpha_1\alpha_2,\\
 a_3&=-\sum \alpha_1\alpha_2\alpha_3,\\
 &\ \vdots\\
 a_n&=\pm \alpha_1\alpha_2\cdots\alpha_n.
\end{aligned}
\tag{3}
\end{equation}

The goal of our considerations is the proof of the principal theorem \gesperrt{that all symmetric functions of the $\alpha$'s can be expressed rationally by the fundamental functions.}

\begin{center}
{\large \S.\ 42.}

\vspace{0.5em}
{\large\gesperrt{The Power Sums.}}
\end{center}

We occupy ourselves first with another special kind of symmetric functions, the \gesperrt{power sums}. Namely, if $\nu$ denotes any positive integral exponent, then
\begin{equation}
 s_\nu=\alpha_1^\nu+\alpha_2^\nu+\cdots+\alpha_n^\nu
\tag{1}
\end{equation}
obviously belongs to the symmetric functions, and $s_\nu$ is called the $\nu$th power sum. We shall derive formulae according to which the power sums are expressible by the fundamental functions.

In what follows, whenever $\varphi(x)$ is any function of $x$, we always denote by
\[
 S[\varphi(\alpha)]
\]
%CONT>>
